% Part: incompleteness
% Chapter: representability-in-q
% Section: prim-rec

\documentclass[../../../include/open-logic-section]{subfiles}

\begin{document}

\olfileid{inc}{req}{pri}
\olsection{شبیه‌سازی بازگشت اولیه}

اکنون می‌توانیم نشان دهیم تعریف با بازگشت اولیه را می‌توان با استفاده
از تابع بتا به‌وسیلهٔ کمینه‌سازی منظم «شبیه‌سازی» کرد. فرض کنید
$f(\vec x)$ و~$g(\vec x,y,z)$ را داریم. آنگاه تابع~$h(\vec x,y)$ که
با بازگشت اولیه از~$f$ و~$g$ تعریف می‌شود چنین است:
\begin{align*}
h(\vec x, 0) & =  f(\vec x) \\
h(\vec x, y+1) & =  g(\vec x, y, h(\vec x, y)).
\end{align*}
باید نشان دهیم می‌توان~$h$ را فقط با ترکیب و کمینه‌سازی منظم از~$f$
و~$g$ تعریف کرد و در این راه توابع پایه و توابعی را به کار برد که
خود با ترکیب و کمینه‌سازی منظم از آن‌ها تعریف شده‌اند (مانند~$\beta$).

\begin{lem}
\ollabel{lem:prim-rec}
اگر بتوان~$h$ را با استفاده از بازگشت اولیه از~$f$ و~$g$ تعریف کرد،
می‌توان آن را از~$f$، $g$ و توابع~$\Zero$، $\Succ$، $\Proj{n}{i}$،
$\Add$، $\Mult$ و~$\Char{=}$ با استفاده از ترکیب و کمینه‌سازی منظم
تعریف کرد.
\end{lem}

\begin{proof}
نخست تابع کمکی~$\hat h(\vec x,y)$ را تعریف کنید که کوچک‌ترین عدد~$d$
را بازمی‌گرداند به‌گونه‌ای که~$d$ دنباله‌ای را رمزگذاری کند که شرایط
زیر را برآورده می‌سازد:
\begin{enumerate}
\item $(d)_0=f(\vec x)$، و
\item برای هر~$i<y$، $(d)_{i+1}=g(\vec x,i,(d)_i)$،
\end{enumerate}
که اینجا $(d)_i$ اختصاری برای~$\beta(d,i)$ است. به بیان دیگر،
$\hat h$ دنبالهٔ
$\tuple{h(\vec x,0),h(\vec x,1),\dots,h(\vec x,y)}$ را
بازمی‌گرداند. می‌توانیم~$\hat h$ را چنین بنویسیم:
\[
\hat h(\vec x, y) = \umin{d}{(\beta(d,0) = f(\vec x) \land \bforall{i <
  y}{\beta(d,i+1) = g(\vec x, i,\beta(d,i)})}.
\]
توجه کنید اینجا هیچ بازگشت اولیه‌ای لازم نیست و فقط کمینه‌سازی به کار
می‌رود. به سبب لم تابع بتا، \olref[bet]{lem:beta}، تابعی که کمینه
می‌کنیم منظم است.

اما اکنون داریم:
\[
h(\vec x, y) = \beta(\hat h(\vec x, y), y),
\]
پس~$h$ را می‌توان تنها با ترکیب و کمینه‌سازی منظم از توابع پایه تعریف
کرد.
\end{proof}

\end{document}
